% !TeX root = ../main.tex

\newlecture[Random Variables \& Distributions][Sep 17, 2026]
\begin{definition}[Density Function][3.7]
    A r.v. has a density function
    \[f: \mathbb{R} \rightarrow \mathbb{R}, \quad \text{if for all }x \in \mathbb{R} \]
    \[F(x) = \int_{-\infty}^{x}f(y)dy\]
    we call $f()$ the density function.
\end{definition}

\begin{definition}[Joint Distribution][3.8]
    We define the joint distribution function of a random vector $\overrightarrow{X} = (X_1, \dots, X_n)$ by
    \[F(x_1, \dots, x_n) = p(X_1 \leq x_1, \dots, X_n \leq x_n)\]
    using the notation $p(A, B) \overset{def}{=}p(A\bigcap B)$
\end{definition}

\begin{theorem}[][3.8.1]
    A joint distribution $F: \mathbb{R}^n \rightarrow \mathbb{R} $ satisfies the following properties:
    \begin{enumerate}[label = \roman*)]
        \item if $(x_1, \dots, x_n)$ and $x_1', \dots x_n'$ satisfy $x_i \leq x_i'$ for all $i \in [n]$.
        \[F(x_1, \dots x_n) \leq F(x_1', \dots x_n')\]
        \item $\lim_{x_i \to -\infty} F(x_1, \dots, x_n) = 0$ for all $i \in \left\{1, \dots, n\right\}$ and 
        $\lim_{x_1 \to \infty, \dots x_n \to \infty}  F (x_1, \dots, x_n) = 1$
        \item $\lim_{h \to 0^+}F(x_1+h, \dots, x_n) = \cdots = \lim_{h \to 0^+}F(x_1, \dots x_n+h)$
        \[=F(x_1, \dots, x_n)\]
    \end{enumerate}
    Proof are anagolous to the proof under univariate case.
\end{theorem}

\begin{definition}[The Joint Density][3.9]
    A random vector has a joint density function:
    \[f: \mathbb{R}^n \to \mathbb{R}, \quad \text{if }\forall(x_1, \dots, x_n)\in \mathbb{R}^n\]
    \[F(x_1, \dots x_n) = \int_{-\infty}^{x_1}\cdots \int_{-\infty}^{x_n}f(y_1, \dots, y_n)\, dy_n\dots \,dy_1\]
\end{definition}

\begin{theorem}[Independence of Random Variables][3.10]
    Random variables $X_1, \dots, X_n$ with CDFs $F_1, \dots, F_n$
    are independent if and only if
    \[F(x_1, \dots, x_n) = \prod_{i=1}^nF_i(x_i), \quad \forall(x_1, \dots, x_n) \in \mathbb{R}^n \]
\end{theorem}

\heading{Limits of sets:}
Let $A_i$ be sets for $i\in\{1,2,\dots\}$ and consider $\displaystyle\lim_{n\to\infty}A_n$.
\begin{itemize}
    \item If $A_1 \subseteq A_2 \subseteq \dots \subseteq A_n \subseteq \dots$ (increasing), then $\displaystyle\lim_{n\to\infty}\bigcup_{i=1}^{n}A_i = \bigcup_{i=1}^{\infty}A_i$.
    \item If $A_1 \supseteq A_2 \supseteq \dots \supseteq A_n \supseteq \dots$ (decreasing), then $\displaystyle\lim_{n\to\infty}\bigcap_{i=1}^{n}A_i = \bigcap_{i=1}^{\infty}A_i$.
\end{itemize}
